
\documentclass[12pt,a4paper]{article}
% \usepackage{lua-visual-debug}

\RequirePackage[no-math]{fontspec}
\newfontfamily\FonteTitulo[
    Path      = ./fontes/qualitype/,
    Extension = .otf,
    BoldFont  = *-Bold,
]{QTBengal}

\setlength\parindent{0pt}
% :FACER: facer esto robusto
\renewcommand*{\baselinestretch}{0} % interlineado global

% :FACER: numero e data en monoespaciado
\makeatletter

% NUMERO
\def\NumeroRevista#1{\def\@NumeroRevista{#1}}
\newcommand{\imprimeNumeroRevista}{Num.\@NumeroRevista}

% DATA
\def\DataRevista#1{\def\@DataRevista{#1}}
\newcommand{\imprimeDataRevista}{\@DataRevista}

% 1º ARTICULO
\def\NomePrimeiroArticulo#1{\def\@NomePrimeiroArticulo{#1}}
\newcommand{\imprimeNomePrimeiroArticulo}{\@NomePrimeiroArticulo}

% TITULO REVISTA
\def\imprimeTituloRevista{A voz da Física}

\makeatother

% Pa colorear o primeiro catacter de algo
% \coloreacolocter{color}{texto}
\newcommand{\coloreacaracter}[2]{%
    {\color{#1}\directlua{tex.sprint(unicode.utf8.sub('#2',1,1))}}%
    \directlua{tex.sprint(unicode.utf8.sub('#2',2,-1))}%
}


% Cores
\usepackage{xcolor}
\definecolor{resalte}{HTML}{ff0000}

% TikZ para os gráficos
\usepackage{tikz}
\usetikzlibrary{calc} % para facer calculos cas posicions
\usepackage{tcolorbox}

\usepackage[
    paper          = a4paper, % tamaño do papel
    top            = 0mm,     % marxe superior
    left           = 0mm,     % marxe esquerdo
    right          = 0mm,     % marxe dereito
    bottom         = 0mm,     % marxe inferior
    nomarginpar    = true,    % desactivar notas no marxe (tamaño 0)
    nohead         = true,    % desactivar o encabezado
    nofoot         = true,    % desactivar o pe de paxina
    % showframe,  % mostra estructura
    % showcrop,   % mostrar tamaño útil (ver opcions layout=xxx)
    % pass        % desactivar todo excepto showframe
]{geometry}

% \usepackage{graphicx}


% En principio lo que buscamos es una portada llamativa pero no excesivamente
% cargada, con una foto en medio pero rodeada con márgenes, que contengan el
% título de la revista, la fecha (mes y año), el número de la revista, el logo de
% la revista, el logo de la USC, y el nombre del primer artículo. La tipografía y
% márgenes se dejará a la imaginación de los autores, así como la colocación de
% cada objeto en la portada. El color dependerá de la foto, si queréis podéis
% hacer alguna prueba.

\NumeroRevista{250}
\DataRevista{Febreiro do 2025}
\NomePrimeiroArticulo{Superconductividade e Superfluidos}

\begin{document}
% Titulo da revista: \imprimeTituloRevista \\
% Titulo da revista coloreado: \coloreacaracter{resalte}{\imprimeTituloRevista} \\
% Revista numero: \imprimeNumeroRevista \\
% Data: \imprimeDataRevista \\
% Primeiro articulo: \imprimeNomePrimeiroArticulo
% \imprimecaracteres{Titulo}

\vspace*{3em}

\begin{tcolorbox}[
    colback = resalte,
    size    = minimal,
    top     = 5pt,
    bottom  = 5pt,
    boxrule = 0pt,
    sharp corners,
]
\end{tcolorbox}

% TITULO E ASOCIADOS
\begin{tcolorbox}[
    colback = white,
    boxrule = 0pt,
    sharp corners,
]
    {%
        \FonteTitulo\fontsize{40pt}{1pt}\selectfont%
        \coloreacaracter{resalte}{\imprimeTituloRevista}%
    }%
    \hfill
    \begin{minipage}[b]{6cm}
        Texto de arriba \\
        Texto abaixo
    \end{minipage}
\end{tcolorbox}

\begin{tcolorbox}[
    colback = resalte,
    size    = small,
    boxrule = 0pt,
    sharp corners,
]
\color{white}{\texttt{\imprimeNumeroRevista\hfill\imprimeDataRevista}}%
\end{tcolorbox}

% IMAXE DE CADA REVISTA
% :FACER: debe cambiar con cada revista, gardar nome nun macro?
\bgroup
\nointerlineskip
\begin{picture}(\textwidth,\textwidth)%
\put(0,0){\hbox{\includegraphics[width=\textwidth]{./black_hole.jpg}}}
\end{picture}
\egroup

% :FACER: variables para manipular as coordenadas do nodo dende os propios
% artigos
\begin{tikzpicture}[remember picture,overlay]
    \node
    [anchor=south west,inner sep=0pt]
    at ($(current page.west)$)
    {%
        \begin{minipage}{0.5\linewidth}
            \begin{tcolorbox}[
                    boxrule     = 0pt,
                    opacityback = 0.5,
                    standard jigsaw,
                    sharp corners,
                ]
                {\fontsize{30pt}{1pt}\selectfont\imprimeNomePrimeiroArticulo}
            \end{tcolorbox}
        \end{minipage}
    };
\end{tikzpicture}

\end{document}
